\section{The combinatorial setting}
\label{sec:setting}

The earlier automata are stated tiling by tiling because they read geometry. A pentagrid rule speaks of five
edge-neighbours and five vertex-neighbours in cyclic order, a heptagrid rule of seven side-neighbours, a
dodecagrid rule of twelve faces. To state one automaton for all of them we drop down to the level on which the
tilings agree, the cell-adjacency graph.

\subsection{A tiling as a graph}

Let $T$ be a regular tiling, of the hyperbolic plane or of hyperbolic space, with cells that meet face to
face. Its cell-adjacency graph $G$ has one vertex for each cell of $T$ and one edge for each pair of cells that
share a facet. Two cells are neighbours exactly when they share a facet. The degree of a vertex of $G$ is the
number of facets of a cell, which is $p$ for the planar tiling $\{p,q\}$ and the facet count of the cell for a
honeycomb \cite{coxeter1973}. We will say cell for a vertex of $G$ and neighbour for an adjacent vertex.

The automaton of Section~\ref{sec:automaton} is defined entirely on $G$. Its rules read a cell, the states of
the cells adjacent to it, and a fixed local labelling described below. They never refer to the angles, the
lengths, the curvature, or the dimension of $T$. The same rules therefore apply to the pentagrid, the
heptagrid, the octagrid, the dodecagrid, and every other regular tiling, because each of these is, at the
level the rules see, just a graph with a labelling.

\subsection{The structured initial configuration}

A weakly universal automaton is allowed an infinite initial configuration, provided it is ultimately periodic,
a finite pattern repeated outward in a regular way. The railway network is such a configuration. We now say
which information it carries. Each cell is in one of three roles, fixed for all time by the initial
configuration.

\begin{itemize}
  \item A cell may be \emph{empty}, not part of the circuit. Empty cells stay empty, and they form the
  quiescent background.
  \item A cell may be a \emph{track} cell. A track cell carries the identities of its two track-neighbours, the
  two cells along the track through it. These two distinguished neighbours are the only ones the locomotive
  uses when passing through the cell.
  \item A cell may be a \emph{switch} cell. A switch cell carries the identities of three track-neighbours in
  labelled roles, the trunk, the branch $a$, and the branch $b$, together with the kind of switch, fixed,
  flip-flop, or memory.
\end{itemize}

This labelling is the oriented part of the construction. It is laid down once, in the initial configuration,
and never changes. It is what lets a switch in any tiling know which of its neighbours is the trunk without
appeal to a global frame. We discuss the status of this labelling, and how it differs from Margenstern's
rotation invariance, in Section~\ref{sec:comparison}. On top of the fixed role and labelling, each cell carries
a dynamic state that the automaton updates, the state that holds the moving locomotive and the switch settings,
the subject of the next section.

\subsection{Laying the circuit on a tiling}

To run a given register machine on a given tiling, one embeds the machine's track network into the cell graph
of the tiling, choosing which cells are track and which are switches and fixing their labels. A hyperbolic
tiling has exponential room, so a track can be routed wherever it is needed and the parts of the circuit can
be kept apart. Two pieces of track that must cross are handled by the crossing of Section~\ref{sec:railway}, in
the plane by the four-way crossing piece and in three or more dimensions by routing one track around the other
through the extra dimension.

We stress that the embedding is part of the input, the program, and not part of the automaton. The automaton
is the fixed rule of Section~\ref{sec:automaton}. Different computations are different initial configurations
of the one automaton, exactly as different railway layouts run different programs on the one set of physical
track parts.

\subsection{The outward direction}
\label{sec:outward}

One piece of the construction needs a sense of outward, the track-building rule of Section~\ref{sec:strong}
grows a register away from the centre, and tracks are most easily laid along branches that lead outward. This
is supplied by the standard addressing of the tiling, the splitting method and its Fibonacci coordinates, which
assigns every cell a tree depth and a path to the root and is established for the regular tilings
\cite{margenstern2007vol1}. We use it only as a tool, the outward direction at a cell is the direction of
increasing depth, and it is laid down with the initial configuration. We do not develop the addressing here.
The exact, dimension-general form of it that the accompanying code uses, including an exact-integer arithmetic
for the coordinates of every regular tiling, is a separate matter from universality and is reported elsewhere
\cite{pollard2026vibe}.
